\scanpage{35}
\mspage{31}
% archive sheet number: 34
% source lines 1--4
Pour $M_c$ donné, et une partie $\Phi\subset M$, on dira que $\Phi$ est \uline{localement finie} si \deleted{$\forall X\in M_c,$}
\[
\forall X\in M_c,\quad
\text{\deleted{posant}\inserted{l'ens.}}\quad
\Phi_X=\{Y\in\Phi\mid Y\mathrel{\overline{\parallel}}X\}
\quad\text{\inserted{est fini} \deleted{on a $\Phi\unsure{\setminus|,}\Phi_X$ fini.}}
\]
% source line 5, closely written above the next paragraph
Les parties loc finies de $M$ forment un \uline{antifiltre}.

% source lines 6--9
L'ensemble \deleted{des} $\unsure{\mathcal F|\mathscr F|\mathfrak F}_{l.f}(M)$ des figures loc. finies de $M$ forme un \inserted{$M$-}atelier, à condition pourtant de supposer \deleted{que} ceci (qui assure que $M\subset\mathcal F_{l.f}(M)$)

% source lines 10--12
\uline{Mag cpt 1}\quad c) \deleted{$\forall\unsure{X|Y}\in M_c,$} $X\in M$, $\widetilde X$ est loc fin pour $M_c$ i.e. $\forall Y\in M_c$, l'ensemble \deleted{des}
\[
\{X'\in\widetilde X\mid X'\mathrel{\overline{\parallel}}Y\}\quad\text{est fini.}
\]

% source lines 13--14
Si $\Phi$ est une préfigure loc finie, alors $\overline\Phi^{(\trianglelefteq)}$ est encore loc finie.

% source lines 15--18
L'axiome Mag at cpt signifie que $\mathcal F\subset\mathcal F_{\mathrm{loc\,fin.}}$. Ainsi $\mathcal F_{\mathrm{loc\,fin}}$ est le plus grand atelier \unsure{adhérant|dérivant|attenant} à $M$, qui soit compatible à la compactologie.

% source lines 19--24
Inversement, donnons nous $(M,\trianglelefteq,\overset{\circ}{\ll},\mathrel{\mid\circ\mid},\mathcal F)$. On va construire une \deleted{une} compactologie la plus grande possible, pour laquelle $M_c,\mathcal F$ soient compatibles, en posant
\[
\begin{aligned}
M_c=\{X\in M\mid {}&\deleted{\unsure{\exists}}\forall F\in\mathcal F,\ \text{l'ens }\{Y\in\widetilde F\mid X\mathrel{\overline{\parallel}}Y\}\\
&\text{est fini}\}.
\end{aligned}
\]
